\subsection{离散群的逆紧且自由的作用}

定理 1. —— 设 G 为从右边连续作用于拓扑空间 E 的离散群。假设 E 的每个点都有邻域 U，使得对 G 的每个异于单位元的元素 g，有 U ∩ (U · g) = ∅。则典范映射 p: E → E/G 使 E 成为 E/G 的以 G 为群的主覆叠。

依假设，G 在 E 上的作用是自由的，故只需证明映射 p 是平展的（I, 第 98 页, 例 1）。设 x 为 E 的点，U 为 x 的开邻域，使得对 G 的每个异于单位元的元素 g 有 \(U \cap U \cdot g = \varnothing\) 。映射 p 是开的（TG, III, 第 10 页, 引理 2），且诱导从 U 到 \(p(\mathrm{U})\) 的单射、连续且开的映射，因而是同胚，这就证明了映射 p 是平展的。

例. —— 1) 设 n 为整数 \(\geqslant 0\) ，\(\mathbf{P}_{n}(\mathbf{R})\) 为 n 维射影空间（TG, VI, 第 13 页），\(S_{n}\) 为 \(R^{n+1}\) 中的单位球面（TG, VI, 第 9 页）。从 \(S_{n}\) 到 \(\mathbf{P}_{n}(\mathbf{R})\) 上的典范映射使 \(S_{n}\) 成为以 \(\{1,-1\}\) 为群的主覆叠，该群以位似作用；轨道是由直径上相对的点组成的偶。

\begin{enumerate}
\def\labelenumi{\arabic{enumi})}
\setcounter{enumi}{1}
\item
  每个次数为 2 的覆叠都具有唯一的以 Z/2Z 为群的主覆叠结构。
\item
  设 X 为 R 上有限维的局部分离微分流形，\(\widetilde{X}\) 为 X 的切丛的定向流形（VAR, R, 10.2.4）。空间 \(\widetilde{X}\) 连同它到 X 上的典范投影，是以 \(\{1,-1\}\) 为群的 X 的主覆叠。
\end{enumerate}

推论 1. —— 设 G 为从右边连续作用于拓扑空间 E 的离散拓扑群。下列条件等价：

\begin{enumerate}
\def\labelenumi{(\roman{enumi})}
\item
  群 G 逆紧且自由地作用于 E；
\item
  空间 E/G 是豪斯多夫的，且 E 是以 E/G 为基、以 G 为群的主覆叠。
\end{enumerate}

此外，在这些条件之下，空间 E 是豪斯多夫的。

设条件 (i) 成立。则空间 E 与 E/G 都是分离的（TG, III, 第 29 页, 命题 3），且对每个 \(x \in \mathrm{E}\) ，存在 \(x\) 在 E 中的开邻域 U，使得 \(\mathrm{U} \cap (\mathrm{U} \cdot g) = \varnothing\) 对 G 的每个异于单位元的元素 \(g\) 成立（TG, III, 第 32 页, 命题 8）。由定理 1，条件 (ii) 随之成立。

蕴涵 (ii)⇒(i) 由 I, 第 96 页推论 1 得出。

推论 2. —— 设 G 为拓扑群，H 为 G 的离散子群。让 H 以右平移作用于 G。则从 G 到模 H 的左陪集空间 G/H 中的典范映射赋予 G 一个 G/H 的以 H 为群的主覆叠结构。

设 V 为 G 的单位元 \(e\) 的邻域，使得 \(\mathrm{H} \cap \mathrm{V} = \{e\}\) 。存在 \(e\) 在 G 中的开邻域 U，使得 \(\mathrm{U}^{-1} \cdot \mathrm{U} \subset \mathrm{V}\) （TG, III, 第 3 页），从而 \(\mathrm{U} \cap \mathrm{U} \cdot h = \varnothing\) 对所有 \(h \in \mathrm{H}\) （\(h \neq e\) ）成立。因此，I, 第 100 页的推论 2 由定理 1 得出。

例 4. —— 赋予从 R 到 T = R/Z 的典范映射（TG, V, 第 2 页），R 是 R/Z 的以 Z 为群的主覆叠。

注 1. —— 设 G 为分离拓扑群，K 为 G 的紧子群，H 为 G 的离散子群。群 G 从右边连续且逆紧地作用于模 K 的右陪集空间 K\G（TG, III, 第 30 页, 推论）。由于子群 H 在 G 中是闭的，它也逆紧地作用于 K\G（TG, III, 第 27 页, 例 1）。此外，若 \(H \cap gKg^{-1} = \{e\}\) 对所有 \(g \in G\) 成立，则群 H 自由地作用于 K\G。于是由推论 1，空间 K\G 是 K\G/H 的以 H 为群的主覆叠。

推论 3. —— 设 G 与 G' 为拓扑群，\(\varphi\colon G \to G'\) 为连续、开且满的同态。假设 \(\varphi\) 的核 H 是离散的。则对于 H 以右平移在 G 上的作用，\(\varphi\) 使 G 成为 G' 的以 H 为群的主覆叠。

此外，若群 G 连通，则 H 包含于 G 的中心。

同态 \(\varphi\) 诱导从 G/H 到 G' 上的拓扑群同构（TG, III, 第 16 页, 命题 24），由此并利用推论 2 即得第一个断言。设群 G 连通。对每个 \(h \in \mathrm{H}\) ，由 \(g \mapsto ghg^{-1}\) 定义的从 G 到 H 的连续映射是常值的，取值为 \(h\) 。故群 H 包含于 G 的中心。

注. —— 2) 设 G 与 G' 为拓扑群，\(\varphi\colon\mathrm{G}\to\mathrm{G}'\) 为连续且开的同态。若群 G' 连通，则同态 \(\varphi\) 是满射。事实上，\(\varphi(\mathrm{G})\) 是 \(\mathrm{G}'\) 的开子群，从而是闭子群，于是等于 \(\mathrm{G}'\)。

\begin{enumerate}
\def\labelenumi{\arabic{enumi})}
\setcounter{enumi}{2}
\tightlist
\item
  设 G 为无穷可数的局部紧群，G' 为基础空间为贝尔空间的分离拓扑群。每个从 G 到 G' 上的连续满同态都是开的（TG, IX, 第 56 页, 推论及 TG, III, 第 16 页, 命题 24）。
\end{enumerate}

例. —— 5) 对每个整数 \(n > 0\)，以 \(\mu_{n}\) 表示 C 中 n 次单位根组成的群（A, V, 第 75 页）。映射 \(z \mapsto z^{n}\) 使 \(C^{*}\) 成为 \(C^{*}\) 的以 \(\mu_{n}\) 为群的主覆叠，该群在 \(C^{*}\) 中以乘法作用。

\begin{enumerate}
\def\labelenumi{\arabic{enumi})}
\setcounter{enumi}{5}
\tightlist
\item
  赋予映射 \(z \mapsto e^{2\pi iz}\) （TG, VIII, 第 8 页, 注），空间 C 是 \(C^{*}\) 的以 Z 为群的主覆叠。映射 \(x \mapsto e^{2\pi ix} = \mathbf{e}(x)\) （从 R 到 \(S_{1}\) ）使 R 成为 \(S_{1}\) 的以 Z 为群的主覆叠。
\end{enumerate}

